\appendix
\onecolumngrid

\section{Notation}
\label{app:notation}

Table~\ref{tab:symbols} defines every symbol of this paper, by group, with the plain-text names under
which the programs, the figures and the data files write it. Two entries are taken from the master
registries of the symbols of PROJECT spacetimes and keep exactly the meaning that they have in the
rest of the project, followed by a note on their use here: the index $j$ (registered in the project as
$i,j$, the index of the space-like coordinates; here only $j$ is an index, and $i$ is the imaginary
unit) and the number $D$ of space dimensions. All other entries are new. Every new symbol carries the
suffix LL (for \emph{Lorentz super-loop}), written as a subscript ($Q_{\rm LL}$) or, where the lower
position holds an index, as a superscript ($\sigma^{\rm LL}_j$), except the following, which cannot be
read as anything else: the bi-octonions and their parts in bold upright type ($\mathbf x$, $\mathbf
e_p$, $\mathbf u$, $\mathbf b$, \dots), the unit matrix $\mathsf 1_4$ and the operators $\mathsf O_1$,
$\mathsf O_2$, $\mathsf O_3$ of Sec.~\ref{sec:nogo}, the number sets and algebras in blackboard type
($\mathbb R$, $\mathbb C$, $\mathbb H$, $\mathbb O$), the groups and Lie algebras ($SL(2,\mathbb C)$,
$\mathfrak{so}(3,1)$, \dots), operation signs, and the plain symbols $i$ (the imaginary unit), $p$ (an
index over the basis of an algebra) and $\iota$ (the isospin index). The letter $\mathsf K$ without
scripts in Theorem~SF1 is any $8\times8$ pair matrix.

The machine checks LL-S1 to LL-S8 (Appendix~\ref{app:checks}) verify that the symbols obey these rules:
no symbol and no plain-text name is registered twice (LL-S1); the taken entries carry the meanings of
the master registries, no new symbol coincides with a symbol of a master registry, and the symbols that
this paper shares with the sibling folder \texttt{Revision\_NoGo} of the repository are declared and
denote the same objects there (LL-S2); every column of every data file is a registered name or a
descriptive word (LL-S3); every document of this revision contains the complete table of symbols
(LL-S4); the symbols of the superfluid, $\mu_{\rm LL}$, $\Lambda_{\rm LL}$, $\Omega_{\rm LL}$,
$\nu^{\rm F}_{\rm LL}$, $\Delta_{\rm LL}$, $\zeta_{\rm LL}$, $\varepsilon_{\rm LL}$, always carry their suffix
(LL-S5); no letter that the project uses with another meaning occurs in a formula without the suffix
LL (LL-S6); every letter in bold, sans-serif, calligraphic, blackboard or fraktur type is registered
(LL-S7); and every name with the suffix \texttt{\_LL}, every name of an algebra element, every
coefficient and every pair matrix that occurs in a report or data file is registered (LL-S8). Units:
$\hbar$ and the speed of light are 1; in Secs.~\ref{sec:superfluid} and~\ref{sec:renorm} the momentum
cutoff $\Lambda_{\rm LL}=1$ unless a unit is shown.

{\small\longtabledepth
\begin{longtable}{p{0.24\textwidth}p{0.54\textwidth}>{\raggedright\arraybackslash\footnotesize}p{0.16\textwidth}}
\caption{The symbols of this paper, by group, with the plain-text names under which the programs, the
figures and the data files write them.}\label{tab:symbols}\\
\toprule
symbol & meaning & plain-text names\\
\midrule
\endfirsthead
\toprule
symbol & meaning (continued) & plain-text names\\
\midrule
\endhead
\bottomrule
\endfoot
{{GENRAW:symbols_ll_texrows}}
\end{longtable}
}\longtabledepthrestore

\emph{Plain-text names.} The programs, the figures and the data files write formulas in plain text,
with the names of the last column of Table~\ref{tab:symbols}. The following conventions complete
that column.
\begin{itemize}
\item Elements of the algebras: \texttt{x}, \texttt{y}, \texttt{z}, \texttt{w} are the bi-octonions
$\mathbf x$, $\mathbf y$, $\mathbf z$, $\mathbf w$; \texttt{aH}, \texttt{bS}, \texttt{cH}, \texttt{dS},
\texttt{bS'} are the bi-quaternions $\mathbf a$, $\mathbf b$, $\mathbf c$, $\mathbf d$, $\mathbf b'$ of
the form $\mathbf a+\mathbf b\,\mathbf e_4$ (the letter H recalls that the element lies in $\mathbb
C\otimes\mathbb H$, the letter S that it multiplies $\mathbf e_4$ to give a spinor charge);
\texttt{uL}, \texttt{uL'}, \texttt{uL1}, \texttt{uL2} are $\mathbf u$, $\mathbf u'$, $\mathbf u_1$,
$\mathbf u_2$ (Lorentz part); \texttt{tI}, \texttt{tI'} are $\mathbf t$, $\mathbf t'$ (isospin);
\texttt{gL}, \texttt{hL} are $\mathbf g$, $\mathbf h$ (elements of $SL(2,\mathbb C)$); \texttt{vI} is
$\mathbf v$ (a unit quaternion, an isospin rotation); \texttt{rP}, \texttt{rP'} are $\mathbf r$,
$\mathbf r'$ (translations). \texttt{e0}, \dots, \texttt{e7} are the units $\mathbf e_0,\dots,\mathbf
e_7$ (and \texttt{1} in \texttt{gL = 1 + i e1 + e2} is $\mathbf e_0$); \texttt{ie1}, \dots, \texttt{ie7}
are $i\mathbf e_1,\dots,i\mathbf e_7$; \texttt{e\_p}, \texttt{e\_j}, \texttt{e\_iota} are $\mathbf e_p$,
$\mathbf e_j$, $\mathbf e_\iota$.
\item Operations: a final \texttt{bar} (\texttt{xbar}, \texttt{bSbar}, \texttt{gLbar}) is the
conjugation $\bar{\mathbf x}$; \texttt{\^{}*} the complex conjugation; \texttt{\^{}dag} the Hermitian
conjugation; \texttt{\^{}T} the transpose; \texttt{\^{}-1} the inverse; \texttt{(x)} the tensor product
$\otimes$; \texttt{o} the composition $\circ$; \texttt{[x,y]} the commutator and \texttt{[x,y,z]} the
associator; \texttt{Jac} the Jacobiator; \texttt{L\_x}, \texttt{R\_x} the left and right multiplications
$L^{\rm LL}_{\mathbf x}$, $R^{\rm LL}_{\mathbf x}$, and \texttt{L\_[e1,e4]} the left multiplication by the
commutator $[\mathbf e_1,\mathbf e_4]$; \texttt{cross(a, b)} is the cross product of two complex
three-component vectors; \texttt{|...|} an absolute value or the length of a vector; the dot in
\texttt{kE.(kE+pE)} the Euclidean scalar product; \texttt{x -> y} a map; \texttt{x(y)} a product
or a function value, as the context of the formula shows.
\item Coefficients of the families: \texttt{c\_tw}, \texttt{c\_sl}, \texttt{c\_PP}, \texttt{c\_Ps},
\texttt{c\_sP}, \texttt{c\_TT}, \texttt{c\_Ts}, \texttt{c\_sT} are $c^{\rm tw}_{\rm LL}$, \dots,
$c^{\rm sT}_{\rm LL}$; a complex one is split as \texttt{ctw\_re + i ctw\_im} into its real and imaginary
parts (and so for the others); \texttt{csP} and \texttt{cTT} are real. In checks M-3, M-4, M-7 and M-9
a failing triple or 4-tuple is given by the positions of its elements in the basis of the family:
positions 0, 1, 2 are $\mathbf e_1$, $\mathbf e_2$, $\mathbf e_3$ and 3, 4, 5 are $i\mathbf e_1$,
$i\mathbf e_2$, $i\mathbf e_3$ (the Lorentz part); in $\mathcal F^{\rm B}_{\rm LL}$ positions 6 to 9 are the
translations $\mathbf e_0$, $i\mathbf e_1$, $i\mathbf e_2$, $i\mathbf e_3$, and in $\mathcal F^{\rm C}_{\rm
LL}$ positions 6 to 8 are the isospin elements $\mathbf e_1$, $\mathbf e_2$, $\mathbf e_3$; the last eight
positions are the spinor charges $\mathbf e_4,\dots,\mathbf e_7$, $i\mathbf e_4,\dots,i\mathbf e_7$.
\item The superfluid: \texttt{dvec} is $\vec d_{\rm LL}$ and \texttt{d\_iota} its components;
\texttt{eps2} is $\mathsf E_{\rm LL}$, \texttt{eps3} the Levi-Civita symbol; \texttt{E\^{}(+-)} stands
for $E^+_{\rm LL}$ and $E^-_{\rm LL}$, and \texttt{E\^{}(+-)(Delta)} for these energies with the gap
\texttt{Delta} in place of $\Delta_{\rm LL}$; \texttt{Delta\_(+-)} are the gaps $\Delta^\pm_{\rm LL}$;
\texttt{Etilde\^{}(+-)} are $\tilde E^\pm_{\rm LL}$; \texttt{eps} is $\varepsilon_{\rm LL}$, \texttt{q} the
momentum $\vec q_{\rm LL}$ or its length; \texttt{Q\_LL(Kch)} is the electric charge of the two spinor
charges that the pair operator $\Psi^T_{\rm LL}\mathsf K\Psi_{\rm LL}$ annihilates; \texttt{0+} and
\texttt{0-} are spin zero with positive and negative parity; \texttt{spin(spin+1)} in check Q-3 is the
spin quantum number times itself plus one, the value 0, 2 or $\tfrac34$ of the Casimir for the spins 0,
1 and $\tfrac12$.
\item Abbreviations: \texttt{BdG} is Bogoliubov-de Gennes, \texttt{BCS} Bardeen-Cooper-Schrieffer;
\texttt{M\_2(C)} the complex $2\times2$ matrices; \texttt{R\^{}8}, \texttt{C\^{}3}, \texttt{C\^{}8} the
real or complex vector spaces of that dimension; \texttt{Im H}, \texttt{Im O} and \texttt{C (x) H} are
${\rm Im}\,\mathbb H$, ${\rm Im}\,\mathbb O$ and $\mathbb C\otimes\mathbb H$; the groups and algebras
(\texttt{SL(2,C)}, \texttt{so(3,1)}, \texttt{g2(C)}, \texttt{SU(2)\_v}, \texttt{U(1)\_F\_LL},
\texttt{Z2}, \dots) are written as in the text; \texttt{4x4} means four by four; ``exact'' in a check
means integer or rational arithmetic, ``double precision'' floating-point arithmetic with 53-bit
mantissas.
\end{itemize}
In the comments of the source code the quaternion units $\mathbf e_1$, $\mathbf e_2$, $\mathbf e_3$ are
also called \texttt{i\_q}, \texttt{j\_q}, \texttt{k\_q}, and the unit $\mathbf e_4$ of the Cayley-Dickson
construction \texttt{l}; these names do not occur in the reports, the data files or the figures.

\section{The machine checks}
\label{app:checks}

The program \texttt{superloop} (Appendix~\ref{app:data}) verifies {{GEN:nchecks_rust}} statements, the
script \texttt{malcev\_groebner.py} the {{GEN:nchecks_groebner}} statements G-1 to G-12 of the
classification; the script \texttt{check\_symbols.py} verifies the eight statements LL-S1 to LL-S8 on the
symbols (Appendix~\ref{app:notation}), and \texttt{check\_self\_contained.py} the statement LL-H1, that
no document of this revision sends its reader to another document for an equation, a symbol or a
definition. Each check prints one line: its label, the statement, and, after two dashes, the numbers
that decide it. A statement marked ``exact'' is decided in exact integer or rational arithmetic, one
marked ``double precision'' with a tolerance far above the rounding errors and far below the effect
that it tests. Table~\ref{tab:checks} reproduces these lines as the programs print them, in plain text
(Appendix~\ref{app:notation} explains the names).

{\footnotesize\longtabledepth
\begin{longtable}{p{0.08\textwidth}p{0.80\textwidth}p{0.06\textwidth}}
\caption{All machine checks of this paper, as the programs print them.}\label{tab:checks}\\
\toprule
label & statement and decisive numbers & result\\
\midrule
\endfirsthead
\toprule
label & statement and decisive numbers (continued) & result\\
\midrule
\endhead
\bottomrule
\endfoot
{{GENRAW:checks_rows_tex}}
\end{longtable}
}\longtabledepthrestore

\section{Programs, data, and how to reproduce every number}
\label{app:data}

\subsection{Programs}

All computations except the Gr\"obner bases are done by the Rust program \texttt{superloop} in the
folder \texttt{code/superloop-rs} of \texttt{Revision\_SuperLoop} ({{GEN:n_rust_files}} source files,
{{GEN:n_rust_lines}} lines); the Gr\"obner bases are computed by the Python script
\texttt{code/scripts/malcev\_groebner.py} ({{GEN:n_groebner_lines}} lines) with Buchberger's algorithm
over the rationals, as implemented in the computer-algebra library SymPy. The figures are drawn by the
Python script \texttt{code/scripts/make\_figures.py} from the data files, and the documents are
assembled by Python scripts that insert every number from the data files, so that no computed number
is typed by hand.

The program has nine modules. \texttt{bioct}: quaternions with complex coefficients and the
Cayley-Dickson product of Eq.~(\ref{eq:cd}) (and, with the opposite sign of the term $\bar{\mathbf
d}\mathbf b$, the split octonions), the conjugations, the norm $N_{\rm LL}$, the commutator, the
associator, the Jacobiator and the linearized Malcev expression~(\ref{eq:malcev}), the left and right
multiplications as real $16\times16$ integer matrices, and a deterministic pseudo-random generator
(SplitMix64) for the random tests; the coefficients are exact Gaussian integers (real and imaginary
parts are 128-bit integers, and an overflow stops the program) or double-precision complex numbers.
\texttt{exact}: exact linear algebra on integer matrices---spans and ranks over the rationals by
fraction-free Gaussian elimination (each row divided by the greatest common divisor of its entries),
Lie closures (commutators are added until the exact span stops growing), the signature of a symmetric
matrix by symmetric elimination in exact rational arithmetic, and ranks modulo the prime $2^{61}-1$ for
the large systems (a rank modulo a prime can only be smaller than the rank over the rationals, so the
dimension of a solution space computed from it is an upper bound; each such bound is matched by the
same number of explicitly verified solutions). \texttt{algebra} (command \texttt{algebra}: checks A, L,
Q and C), \texttt{malcev} (command \texttt{malcev}: the brackets of Eq.~(\ref{eq:families}) with
integer or Gaussian-integer coefficients, the Jacobi identity on all basis triples and the linearized
Malcev identity on all basis 4-tuples, $14^4=38416$ for $\mathcal F^{\rm A}_{\rm LL}$, $18^4=104976$ for
$\mathcal F^{\rm B}_{\rm LL}$ and $17^4=83521$ for $\mathcal F^{\rm C}_{\rm LL}$; checks M),
\texttt{manifolds} (command \texttt{manifolds}: checks N), \texttt{superfluid} (command
\texttt{superfluid}: the Dirac matrices, the pair channels, the Bogoliubov-de Gennes matrices, the
grand potential and the gap equation; checks S), \texttt{renorm} (command \texttt{renorm}: power
counting, the contact coupling, the bubble, the invariants; checks R), \texttt{checks} (the reports) and
\texttt{main} (the commands, and a guard that lets the program run only in the folder
\texttt{Revision\_SuperLoop}, into whose folder \texttt{data} it writes). The library crate
\texttt{code/spacetimes-rs} of the repository supplies complex numbers, the diagonalization of complex
Hermitian matrices, Gauss-Legendre nodes and weights, Brent's root finder and the writing of data files.

\emph{Numerical methods.} The statements of the checks A, L, C, M, N, Q-3, Q-4, Q-6, R-1, R-2, R-5
and R-6 are decided in exact integer or rational arithmetic. The other checks use double precision:
the eigenvalues of checks Q-1, Q-2 and Q-5; the structural checks of the superfluid, S-0, S-1, S-2,
S-7, S-8 and S-9, whose matrices have the entries 0, $\pm1$, $\pm i$ and $\pm\tfrac12$, so that their
products are exact or rounded at the level $10^{-15}$, compared with tolerances between $10^{-14}$ and
$10^{-10}$ (a rank is counted as the number of eigenvalues of the Gram matrix of the vectors that exceed
$10^{-10}$ times the largest); the spectra of check S-3, from the diagonalization of the $16\times16$
Bogoliubov-de Gennes matrices at 60 random points, with random momenta, masses, chemical potentials and
order parameters; and the integrals of checks S-4 to S-6, S-10, R-3 and R-4. The momentum integrals of Sec.~\ref{sec:superfluid} use Gauss-Legendre rules of 40 points on panels whose
end points are 0, the cutoff, and 0.5, 0.9, 0.99, 1, 1.01, 1.1 and 1.5 times the Fermi momentum $k^{\rm
F}_{\rm LL}$ (where these lie inside the cutoff); those of check R-3 use rules of 24 points with end points
at 0.5, 0.9, 0.99, 1, 1.01, 1.1 and 2 times $k^{\rm F}_{\rm LL}$ and further panels whose lengths double
up to the cutoff. The gap equation~(\ref{eq:gapeq}) is solved for $\Delta_{\rm LL}$ by Brent's method
between $10^{-9}$ and 20 with the tolerance $10^{-13}$; the minimum of $\Omega_{\rm LL}$ over
$\lvert\vec d_{\rm LL}\rvert$ is found by golden-section search (120 or 150 steps) on an interval from 0 to
three times the gap. The slope of check S-5 is taken from the two weakest couplings of the grid of 40
couplings, the slope of check R-4 from the two largest cutoffs. The bubble~(\ref{eq:bubble}) is
integrated in four-dimensional spherical coordinates: the length of the loop momentum on panels whose
lengths grow by the factor 1.5 from 0.25 up to the cutoff, and the angle between the loop momentum and
the external momentum from 0 to $\pi$ with the weight $4\pi$ times the square of its sine (the volume of
the three-sphere of directions), each with Gauss-Legendre rules of 48 points.

\subsection{Commands}

From the folder \texttt{Revision\_SuperLoop} of the repository (a Rust toolchain for the edition 2021,
Python 3 with SymPy, NumPy and Matplotlib, and for the documents a TeX distribution with REVTeX 4.2,
pandoc and XeLaTeX):
\begin{verbatim}
cargo run --release --manifest-path code/superloop-rs/Cargo.toml -- all
python3 code/scripts/malcev_groebner.py
python3 code/scripts/make_figures.py
bash code/scripts/check_all.sh
\end{verbatim}
The first command writes the data files of the program and the reports of its checks into the folder
\texttt{data} (in about one second); the commands \texttt{algebra}, \texttt{malcev},
\texttt{manifolds}, \texttt{superfluid} and \texttt{renorm} in place of \texttt{all} run one part. The
second computes the Gr\"obner bases and the checks G-1 to G-12 (a few seconds). The third draws the
seven figures. The fourth runs all three, then the checks of the symbols and of self-containedness,
writes the fragments of the documents, builds this paper and the provenance documents, and ends with
one verdict.

\subsection{Data files}

{{GENRAW:data_files_tex}}
Every data file states its content and units in its header lines (\texttt{\# content:}), and every
column name is either a descriptive word or one of the names listed in the last column of
Table~\ref{tab:symbols}. The reports \texttt{sl\_algebra\_checks.txt}, \texttt{sl\_malcev\_checks.txt},
\texttt{sl\_manifolds\_checks.txt}, \texttt{sl\_superfluid\_checks.txt}, \texttt{sl\_renorm\_checks.txt},
\texttt{sl\_groebner\_checks.txt}, \texttt{sl\_symbol\_checks.txt} and
\texttt{sl\_self\_contained\_checks.txt} contain the lines of Table~\ref{tab:checks}, and
\texttt{sl\_groebner\_bases.txt} the reduced Gr\"obner bases of all six ideals of
Table~\ref{tab:groebner}.

\subsection{Checksums}

SHA-256 checksums of the data files, the figures and the programs from which this paper was built:
{\footnotesize
\begin{verbatim}
{{GENRAW:data_sha256}}
{{GENRAW:figure_sha256}}
{{GENRAW:code_sha256}}
\end{verbatim}
}

\twocolumngrid
